
\documentclass[11pt]{article}
\usepackage[utf8]{inputenc}
\usepackage{amsmath,amssymb,amsfonts}

\title{CNSH Unified Theory Layer}
\author{UID9622 · 龍芯北辰}

\begin{document}

\section{Axioms}

\subsection{Axiom 1: Traceability}
\textbf{Principle:} No execution may occur without an associated trace identifier.

\textbf{Formalization:} $∀E ∈ Executions: ∃DNA(E) ∈ TraceSpace$

\textbf{Application:} All system operations must have unique DNA codes

\subsection{Axiom 2: Immutability}
\textbf{Principle:} Once recorded, trace cannot be modified without detection.

\textbf{Formalization:} $modify(trace) → hash_break ⊥ valid$

\textbf{Application:} Hash chains protect audit logs from tampering

\subsection{Axiom 3: Recoverability}
\textbf{Principle:} System must maintain checkpoints for state recovery.

\textbf{Formalization:} $∀state: ∃checkpoint ∧ checkpoint.valid$

\textbf{Application:} Rollback to safe states on critical failure

\subsection{Axiom 4: Sovereignty}
\textbf{Principle:} Identity cannot be delegated or transferred.

\textbf{Formalization:} $identity(A) ≠ identity(B) ∀A≠B$

\textbf{Application:} Each entity has permanent, unique identity

\subsection{Axiom 5: Causality}
\textbf{Principle:} Execution order defines causal relationships.

\textbf{Formalization:} $t_i < t_j ⟹ causally_independent(e_i, e_j) ∨ depends(e_j, e_i)$

\textbf{Application:} Preventing causality violations in distributed systems

\section{Theorems}

\subsection{Theorem 1.1: Runtime Identity Axiom}
Given: $A system S with semantic continuity property C$

When: $S maintains traceable execution history H$

Then: $S possesses immutable identity I = f(C, H)$

\textbf{Proof:} By definition of runtime identity (Def 1.1), identity I is composed of semantic continuity and traceable history. Given C and H, I is uniquely determined.

\end{document}

\subsection{Theorem 1.2: Traceability Chain}
Given: $Execution trace T = {t_0, t_1, ..., t_n}$

When: $Each t_i has hash h_i = SHA256(t_i || h_{i-1})$

Then: $Chain integrity preserved: ∀ modifications detected$

\textbf{Proof:} Hash chain property ensures cryptographic immutability. Any modification to t_i changes h_i, breaking downstream hashes.

\end{document}

\subsection{Theorem 1.3: Dynamic Digital Root}
Given: $State parameters: N (number), S (sum), R (root), T (time)$

When: $Weight factors α=0.35, β=0.25, γ=0.25, δ=0.15$

Then: $DR = 0.35N + 0.25S + 0.25R + 0.15T ∈ [1, 9]$

\textbf{Proof:} By weighted summation and modulo reduction. Proof by construction (see appendix A).

\end{document}

\subsection{Definition 1.1: Longhun Runtime Identity}
Given: $Semantic continuity C, execution history H, audit log A$

When: $All three properties are well-defined$

Then: $Identity I ≡ ⟨C, H, A⟩ represents immutable runtime entity$

\textbf{Proof:} By definition. Identity is the tuple of three core properties.

\end{document}

\subsection{Axiom 1: Traceability Axiom}
Given: $Any execution E in system S$

When: $E occurs at time t$

Then: $∃ trace identifier DNA(E, t) such that E is permanently recordable$

\textbf{Proof:} Axiomatic - foundational principle

\end{document}

\subsection{Corollary 1.1: Immutability from Traceability}
Given: $Axiom 1 (Traceability) holds$

When: $Trace is protected by cryptographic hash$

Then: $Execution history cannot be modified without detection$

\textbf{Proof:} Follows from hash chain property. See Theorem 1.2.

\end{document}

\subsection{Lemma 1.1: Hash Collision Probability}
Given: $Hash function h: SHA256 with output space 2^256$

When: $Computing collision probability P_c for n executions$

Then: $P_c ≈ n²/2^257 (birthday paradox bound)$

\textbf{Proof:} By probabilistic analysis. Collision bound derived from birthday attack complexity.

\end{document}

\subsection{Theorem 2.1: Semantic Drift Detection}
Given: $Semantic state S_t at time t, S_{t-1} at t-1$

When: $Using Euclidean distance metric in semantic space$

Then: $Drift D(t) = ||S_t - S_{t-1}|| ∈ [0, ∞)$

\textbf{Proof:} By definition of Euclidean metric. Drift measures semantic displacement.

\end{document}

\subsection{Theorem 2.2: Trust Score Composition}
Given: $Integrity I ∈ [0,1], Traceability T ∈ [0,1], Recoverability R ∈ [0,1]$

When: $Trust S = I × T × R$

Then: $S ∈ [0, 1] represents overall system trustworthiness$

\textbf{Proof:} Multiplicative composition ensures all components must be positive.

\end{document}

\subsection{Formula: Five Elements Dynamic Mapping}
Given: $Five elements (木火土金水) as state variables$

When: $Mapping to runtime states (generate, process, archive, secure, flow)$

Then: $Bijection F: {Five Elements} → {Runtime States}$

\textbf{Proof:} By semantic correspondence. Each element maps to distinct execution phase.

\end{document}

\subsection{Theorem 3.1: State Machine Completeness}
Given: $State set Q = {q_0, q_1, ..., q_n}$

When: $Transition function δ: Q × Σ → Q is total$

Then: $All reachable states are covered by some execution path$

\textbf{Proof:} By induction on execution length. Every state reachable in finite steps.

\end{document}

\subsection{Theorem 4.1: Entropy Conservation}
Given: $Information set I with probability distribution p$

When: $Computing Shannon entropy H(I) = -Σ p_i log p_i$

Then: $H(I) ≥ 0, with equality iff one p_i = 1$

\textbf{Proof:} By properties of logarithm and probability distribution.

\end{document}

\subsection{Theorem 5.1: Circuit Breaker Pattern}
Given: $Failure threshold T_f, monitoring window W$

When: $Failure rate exceeds T_f in window W$

Then: $System enters OPEN state, preventing cascading failures$

\textbf{Proof:} By fault tolerance theory. Early termination prevents amplification.

\end{document}

\subsection{Theorem 6.1: Hierarchical Routing}
Given: $Route graph R = (V, E) with weight function w$

When: $Dijkstra's algorithm computes shortest path$

Then: $Optimal routing guaranteed with O(|V|² + |E|) complexity$

\textbf{Proof:} By Dijkstra optimality conditions.

\end{document}

\subsection{Corollary 2.1: Information Flow Bounds}
Given: $Entropy H(I), channel capacity C$

When: $I transmitted through channel with capacity C$

Then: $Transmission possible iff H(I) ≤ C$

\textbf{Proof:} By Shannon's noisy channel coding theorem.

\end{document}

\subsection{Lemma 2.1: Path Existence in DAG}
Given: $Directed acyclic graph G = (V, E)$

When: $Topological ordering exists$

Then: $Any two nodes have at most one path between them (in transitive closure)$

\textbf{Proof:} By acyclicity property. Cycles prevent multiple paths.

\end{document}

\subsection{Theorem 7.1: Rollback Safety}
Given: $System state snapshot S_i at checkpoint i$

When: $Failure occurs at state S_j (j > i)$

Then: $Rollback to S_i preserves all invariants$

\textbf{Proof:} By checkpoint consistency guarantees.

\end{document}
